\cotefile{15}
\scanpage{103}
\mspage{IV.8}

$\psi^*\psi_*(\handwritten{\mathscr{A}}')$ to $\handwritten{\mathscr{A}}'$ corresponding to the identity homomorphism on
$\psi_*(\handwritten{\mathscr{A}}')$. We denote this map by $\sigma_{\handwritten{\mathscr{A}}'}:\psi^*\psi_*(\handwritten{\mathscr{A}}')\handwritten{\longrightarrow}\handwritten{\mathscr{A}}'$. If $\handwritten{\mathscr{A}}$ is a
sheaf of rings on $X$, there is a canonical ring homomorphism from
$\handwritten{\mathscr{A}}$ to $\psi_*\psi^*(\handwritten{\mathscr{A}})$ corresponding to the identity homomorphism on $\psi^*(\handwritten{\mathscr{A}})$.
We denote this map by $\rho_{\handwritten{\mathscr{A}}}:{}\handwritten{\longrightarrow}\psi_*\psi^*(\handwritten{\mathscr{A}})$.

Now suppose $\Psi:(Y,\handwritten{\mathscr{A}}')\handwritten{\longrightarrow}(X,\handwritten{\mathscr{A}})$ is a map of ringed spaces ;
i.e. $\Psi=(\psi,\theta)$ where $\psi$ is a continuous map from $Y$ to $X$ and $\theta$
is a homomorphism of sheaves of rings $\theta:\handwritten{\mathscr{A}}\handwritten{\longrightarrow}\psi_*(\handwritten{\mathscr{A}}')$. Then for
any $\handwritten{\mathscr{A}}'$-module $\handwritten{\mathscr{M}}'$, $\psi_*(\handwritten{\mathscr{M}}')$ can be made into an $\handwritten{\mathscr{A}}$-module by
means of the operation of $\theta$. When $\psi_*(\handwritten{\mathscr{M}}')$ is considered a module
over $\handwritten{\mathscr{A}}$ in this way, it is denoted $\Psi_*(\handwritten{\mathscr{M}}')$. We have the following
canonical functorial homomorphisms [8, 0 ; \S\ 4.2] :

\[
\begin{aligned}
h_* &: \Psi_*(\underline{\mathrm{Hom}}_{\handwritten{\mathscr{A}}'}(\handwritten{\mathscr{M}}',\handwritten{\mathscr{N}}'))\handwritten{\longrightarrow}\underline{\mathrm{Hom}}_{\handwritten{\mathscr{A}}}(\Psi_*(\handwritten{\mathscr{M}}'),\Psi_*(\handwritten{\mathscr{N}}'))\\
t_* &: \Psi_*(\handwritten{\mathscr{M}}')\otimes_{\handwritten{\mathscr{A}}} \Psi_*(\handwritten{\mathscr{N}}')\handwritten{\longrightarrow}\Psi_*(\handwritten{\mathscr{M}}'\otimes_{\handwritten{\mathscr{A}}'}\handwritten{\mathscr{N}}').
\end{aligned}
\]

On the other hand we have the map $\psi^*(\theta):\psi^*(\handwritten{\mathscr{A}})\handwritten{\longrightarrow}\psi^*\psi_*(\handwritten{\mathscr{A}}')$, and
so we have $\sigma_{\handwritten{\mathscr{A}}'}:\psi^*(\theta):\psi^*(\handwritten{\mathscr{A}})\handwritten{\longrightarrow}\handwritten{\mathscr{A}}'$. $\handwritten{\mathscr{A}}'$ is an algebra over $\psi^*(\handwritten{\mathscr{A}})$
by the means of this map, and if $\handwritten{\mathscr{M}}$ is an $\handwritten{\mathscr{A}}$-module, we define
the $\handwritten{\mathscr{A}}'$-module $\Psi^*(\handwritten{\mathscr{M}})$ to be $\psi^*(\handwritten{\mathscr{M}})\otimes_{\psi^*(\handwritten{\mathscr{A}})}\handwritten{\mathscr{A}}'$.

There is a canonical functorial isomorphism [8, 0; 4.33.1]

\[
t^*:\Psi^*(\handwritten{\mathscr{M}})\otimes_{\handwritten{\mathscr{A}}'} \Psi^*(\handwritten{\mathscr{N}})\qquad\Psi^*(\handwritten{\mathscr{M}}\otimes_{\handwritten{\mathscr{A}}}\handwritten{\mathscr{N}}).
\]

It can be shown that there is a natural isomorphism from
$\mathrm{Hom}_{\handwritten{\mathscr{A}}}(\Psi^*(\handwritten{\mathscr{M}}),\handwritten{\mathscr{M}}')$ to $\mathrm{Hom}_{\handwritten{\mathscr{A}}'}(\handwritten{\mathscr{M}},\Psi_*(\handwritten{\mathscr{M}}'))$ for all $\handwritten{\mathscr{A}}$-modules $\handwritten{\mathscr{M}}$
and $\handwritten{\mathscr{A}}'$-modules $\handwritten{\mathscr{M}}'$. By means of this isomorphism we can define
the $\handwritten{\mathscr{A}}$-homomorphism $\rho_{\handwritten{\mathscr{M}}}:\handwritten{\mathscr{M}}\handwritten{\longrightarrow}\Psi_*\Psi^*(\handwritten{\mathscr{M}})$ and the $\handwritten{\mathscr{A}}'$-homomorphism

% archive sheet number: 102

\scanpage{104}

Galois motiviques \unsure{sur|en|de} $\overline{\mathbb F}_p$.
\margin{On \unsure{aura|a donc} $M(q^\infty)\simeq\mathbb Z\times M^{\unsure{0|\circ}}(q^\infty)$, $M^0(q^\infty)=\underset{\longrightarrow}{\lim}\,M^0(q^n)\simeq M^0\otimes_{\mathbb Z}\mathbb Q$}
\deleted{\unsure{Les applications|D'ailleurs les} $M^0(q^{\deleted{\unsure{n|m}}\,\prime})\unsure{-|\longrightarrow} M^0(q)$ induisent un \unsure{isom.}
$M^0(q')\otimes_{\mathbb Z}\mathbb Q\unsure{-|\longrightarrow}M^0(q)\otimes_{\mathbb Z}\mathbb Q$,}
on constate que
$M^0(q^\infty)$ \deleted{qui} s'identifie \unsure{:|à} \deleted{\unsure{au sous-groupe}}
\unsure{à} $M^0\deleted{\unsure{(q)}}\otimes_{\mathbb Z}\mathbb Q$, qui est lui-même isomorphe
au quotient $P/U$, $U$ étant le groupe des
racines de l'1, et $P$ étant le sous-groupe
de $\overline{\mathbb Q}^{*}$ formé des $\lambda$ tels que $\lambda^n\in M^0\deleted{\unsure{(q)}}$
pour $n$ assez grand. \deleted{\unsure{L'injection|La flèche|L'homom.}}

\deleted{\[
M^0(q^{\deleted{\unsure{n}}}=p^n)\unsure{\longrightarrow|-}P/U
\]}

\deleted{s'obtient par}

\deleted{\[
\lambda\unsure{\longrightarrow|\longmapsto}\lambda^{1/n}\bmod U
\]}

\deleted{(où la détermination de $U$ qu'on prend est
sans importance). On a bien $(\lambda^{1/n})^{\deleted{\unsure{mn}}n}=\lambda^{\deleted{\unsure{m}}}\in M^0(p^n)$,
car $M(p^n)^{\unsure{n|1|\wedge}}\unsure{\subset|\in} M(p)$. Trouve-t-on $P^0/U$ tout entier ? $M^{\unsure{0|n|\wedge}}(p)$}

\deleted{\deleted{\unsure{Oui, car}} Oui $\lambda\in\overline{\mathbb Q}^{*}$ tel que $\lambda^n\in M^0(p)$, \deleted{\unsure{donc}}
provient-il de $\mu\in M^0(p^m)$, i.e. \deleted{est-il de}
a-t-on $\lambda^m\in M^0(p^m)$ pour $m$ convenable ??}

\deleted{On peut \deleted{\unsure{supposer, de supposer}} se borner aux $m={}$\unsure{mls}
multiples de $n$, et la question serait si pour
\unsure{tt} $\lambda\in M^0(p)$, $n\deleted{\unsure{,m}}\in\mathbb N^*$ étant donnés, il existe un $h\in\mathbb N^{\unsure{*|k}}$
tel que $\lambda^h\in M^0(p^{hn})$.}

% archive sheet number: 103

\scanpage{105}
\mspage{\handwritten{94}}

be a basis of the group $H^{\unsure{0|O}}(V,O_V[N\tilde H])$. We have:

\[
\handwritten{f_*\left(O_V[N\tilde H]\right)\simeq(f_*O_V)[NH]_{\bar V}.}
\]

From the isomorphism of ringed spaces $(V,f^*O_{\bar V})$ and $(\bar V,O_{\bar V})$
we obtain the isomorphism of the sheaves $O_V[N\tilde H]$ and
$(f_*O_V)[NH]_{\bar V}$.

Let $\bar\alpha,\bar y,\bar y_1,\ldots,\bar y_M$ be the elements over $\bar V$, corresponding
to the elements $\alpha,y,y_1,\ldots,y_M$. Since $\alpha\,y\in O_V[N\tilde H](P)$, then

\[
\handwritten{\bar\alpha\,\bar y\in(f_*O_V)[NH]_{\bar V}(\bar P).}
\]

According to the definition of the number $N$, there
exist $\beta_\nu\in O_{\bar V,\bar P}$ such that

\[
\handwritten{\bar\alpha\,\bar y=\sum_{\nu=1}^{M}\beta_\nu\bar y_\nu}
\]

Hence, if we denote by the same symbol the functions,
corresponding to $\beta_\nu$ at the mapping $f$, we obtain:

\[
\handwritten{\alpha\,y=\sum_{\nu=1}^{M}\beta_\nu  y_\nu}
\]

or, since $y\ne0$ in $P$

\[
\handwritten{\alpha=\sum_{\nu=1}^{M}\beta_\nu\frac{y_\nu}{y}.}\tag*{\handwritten{(24)}}
\]

% archive sheet number: 104

\scanpage{106}

La forme plus faible, ``géométrique'' de la
conjecture ci-dessus dirait que \deleted{\unsure{c'est|aussi}} l'on obtient
bien ainsi, par le \unsure{schémoïde|schémaïde} associé à
$M^{\unsure{0|\circ}}\otimes_{\mathbb Z}\mathbb Q$ (sur lequel $\pi$ opère \ldots), la
\deleted{\unsure{opération|grande}} ``partie compacte'' du \deleted{\unsure{sous}} groupe de Galois
motivique de $\overline{\mathbb F}_p$, qui serait donc isom. à

\[
G\simeq D(M_0\otimes\mathbb Q)\times\mathbb G_m
\]

avec \unsure{l'hom|les hom|l'homom.} inj : $\mathbb G_m\longrightarrow G$ s'identifiant
à $j$, et \deleted{$\unsure{(x,\lambda)}$} $(x,\lambda)\unsure{\longmapsto|-}\lambda^{2\deleted{\unsure{x}}}$ s'identifiant à $\varepsilon$.

En termes de valeurs propres de Frobenius :

\uline{Conjecture faible}. Tout $\lambda\in\overline{\mathbb Q}^{*}$, \deleted{\unsure{de la|tel que} tel qu'il}
% slanted
\margin{\unsure{Résultera|Résulterait} \deleted{\unsure{peut}} \unsure{probablement} de la théorie des multiplications complexes : \unsure{cf|cf.} \uline{\unsure{Tanyama-Shimura|Taniyama-Shimura}}}
\deleted{existe un entier $n>0$ avec $\lambda^n$ satisfaisant aux
conditions, à \unsure{savoir} p} un \uline{entier} algébrique,
\deleted{tel qu'il existe un entier $n$}
tel qu'il existe un entier $i\geq0$, de sorte
que $\lambda=p^{i/2}u$, \deleted{soit un} où $u$ est une \uline{unité}
\uline{absolue} \uline{exc}. $p$, et qu'il existe un entier
$n>0$ tel que $u^n\in\mathbb Q'(i)^*$, et $v_p(u^n)+v_{p'}(u^n)=0$
pour \unsure{ts|tt} couples d'idéaux premiers de $\mathbb Q'(i)$ \unsure{induisant|relevant|au-dessus}
\unsure{sur le|sur|d'un} \unsure{$=$|même} idéal premier de $\mathbb Q'$, Alors il
existe un entier $m>0$ tel que
$\lambda^m=(p^m)^{i/2}u^m$ soit \inserted{une} valeur propre de frobenius \unsure{$f_{q^m}$|géom.|$f_{p^m}$}
dans un motif semi-simple \inserted{de poids $i$} convenable sur $\mathbb F_{\unsure{q^m|p^m}}$

% archive sheet number: 105

\scanpage{107}
\mspage{\handwritten{92}}

$H^{\unsure{0|O}}(V,O_V[nD])$. Let us examine a covering $\{U_\mu\}$ of the
variety $V$, where $U_\mu$ is an addition to the variety of
zeros of the cross-section $\psi_\mu, \mu=\overline{1,R}$.

Since $\psi_\mu$ do not have common zero, then $\{U_\mu\}$
covers all $V$. Some hyperplane cross-section $H_\mu$ of the
variety $\bar V$ becomes the image of the variety of zeros of
the cross-section $\psi_\mu$ at the mapping $f$. The family
$\{U_\mu\}$, where $\bar U_\mu=fU_\mu$, forms an open covering of the
variety $\bar V$. The fibre $[H]_{\bar V}$ is defined for this
covering by the transition function:

\[
\handwritten{\bar R_{\mu_1\mu_2}=\frac{\psi_{\mu_2}}{\psi_{\mu_1}}\quad\text{in } U_{\mu_1}\cap  U_{\mu_2}}
\]

Consider a covering $\{U_{i\mu}\}$ of the variety $V$, where

\[
\handwritten{U_{i\mu}=U^{(i)}\cap  U_\mu}
\]

($\{U^{(i)}\}$ is the covering of $V$, in which $[D]$ was defined
originally). For $\{U_{i\mu}\}$, $[nD]$ is defined by the transition
functions:

\[
\handwritten{{}_1R_{i_1\mu_1,i_2\mu_2}=R^n_{i_1i_2}\quad\text{in }U_{i_1\mu_1}\cap  U_{i_2\mu_2}}
\]

($R_{ij}$ are the transition functions for $[D]$ and $\{U^{(i)}\}$);
$[\tilde H]$ is defined for $\{U_{i\mu}\}$ by the transition functions:

\[
\handwritten{{}_2R_{i_1\mu_1,i_2\mu_2}=\frac{\psi_{\mu_2}}{\psi_{\mu_1}}=\frac{\psi_{\mu_2}^{(i_1)}}{\psi_{\mu_1}^{(i_{\unsure{1|2}})}}\quad\text{in }U_{i_1\mu_1}\cap  U_{i_2\mu_2}}
\]

% archive sheet number: 106

\scanpage{108}

$K$ extension tot. réelle [i.e. si $K$ défini par $H\subset\pi_1(\mathbb Q)$, $H$
ne contient aucun conjugué de $\tau$,
$\tau$ la conjugaison ordinaire]

A) $L$ extension quadratique, $L=K(\unsure{e|x})$, $x^2+ax+b=0$ \quad$(a,b\in K)$.
\inserted{$\sigma$ l'élt non-trivial de $\mathrm{Gal}(L/K)$, alors $\sigma x=\bar x$ pour tt $x$}
\margin{$a=x+\deleted{\unsure{x}}\bar x$, $b=x\deleted{\unsure{x}}\bar x$}
Conditions équivalentes :

\diagram{\begin{tikzcd}[ampersand replacement=\&,row sep=small]
{\text{$L$ tot. imaginaire}}\arrow[d,Leftrightarrow]\\
{\begin{gathered}\text{Pour tt plongmt de $K$ dans $\mathbb R$, }\Delta=\deleted{\unsure{b}}a^2-4b\\
=(x-\deleted{\unsure{x}}\bar x)^2\text{ est }<\deleted{\unsure{\geq}}0,\ \text{i.e. $\Delta$ est}\\
\text{tot. \deleted{positif} \inserted{str.} négatif}\end{gathered}}\arrow[d,Leftrightarrow]\\
{\deleted{\begin{gathered}\text{$L$ est stable par $\unsure{x|t}\mapsto\bar x$, et}\\
(x-\bar x)^2\text{ est tot. positif}\end{gathered}}}
\end{tikzcd}}

\deleted{Alors $\sigma$ est la conj. complexe \unsure{ord.}}

B) Soit $x\in\overline{\mathbb Q}\subset\mathbb C$. Conditions équivalentes

\diagram{\begin{tikzcd}[ampersand replacement=\&,row sep=small]
{1)\ \left\{\begin{aligned}&x+\bar x\ [\text{et }x\bar x\text{ sont}]\ \inserted{\text{conséquence des autres, car }x\bar x=\frac14[(x+\bar x)^2-(x-\bar x)^2]}\ \text{tot. réels}\\
&(x-\bar x)^2\text{ est tot. \deleted{\unsure{positif $>0$|$\geq0$} positif \unsure{$<0$}} \inserted{str.} négatif}\end{aligned}\right.}\arrow[d,Leftrightarrow]\\
{2)\quad\begin{gathered}\deleted{\text{$x$ est quadr.}}\ \exists\text{ un $K$ tot. réel et une}\\
\text{extension $L$ tot. imaginaire de $K$, dont}\\
\text{$x$ soit un générateur.}\end{gathered}}\arrow[d,Leftrightarrow]\\
{2')\quad\deleted{\unsure{\exists}}\,x\in\mathbb Q'(i)-\mathbb Q',\quad\text{i.e.}}
\end{tikzcd}}

En effet, $2)\Rightarrow1)$ en vertu de A), et pour $1)\Rightarrow2)$, il suffit
de prendre $K=\mathbb Q(a,b)$ où $a=x+\bar x$, $b=x\bar x$, et appliquer
encore A).

\uline{Remarque} Cel\unsure{à|a} prouve que les
conditions de 1) sont stables par conjugaison, ce qui
n'était pas évident autrement.

\deleted{3) $x\bar x$ est tot. str. positif
et de même pour $x_\alpha\bar x_\alpha$ pour tt conjugué de $x$}

C) Soit $x\in\overline{\mathbb Q}\subset\mathbb C$. Conditions équivalentes

\diagram{\begin{tikzcd}[ampersand replacement=\&,row sep=small]
{1)\quad x+\bar x\text{ tot. réel},\quad(x-\bar x)^2\text{ tot. négatif}}\arrow[d,Leftrightarrow]\\
{2)\quad\exists\ K,L\text{ comme plus haut},\quad x\in\deleted{K}L}\arrow[d,Leftrightarrow]\\
{3)\quad\unsure{a|x}\in\mathbb Q'(i).}\arrow[d,Rightarrow,bend right=28]\\
{4)\quad x\bar x\text{ tot. positif, \deleted{et de même} \deleted{\unsure{équiv.|égal}} \inserted{ainsi que} }x_\alpha\bar x_\alpha\text{ pour tt conjugué $x_\alpha$ de $x$.}}\arrow[u,Rightarrow,"{?}"',bend left=55]
\end{tikzcd}}

% archive sheet number: 107

\scanpage{109}
\mspage{-10-}

$G_{x\unsure{0|O}}$, \uline{the intersection} $G'\cap T'$ \uline{finite,} \uline{and} $\operatorname{rank}G'+\operatorname{rank}T'=r$. \uline{Further}
$G'$ \uline{and} $T'$ \uline{are uniquely determined as the commutator subgroup and the
identity component of the centre of} $G_{x0}$. (b) \uline{The unipotent elements of}
$G_x$ \uline{are all in} $G'$.

Part (b) follows from (a) because $G_{x0}$ contains the unipotent elements
of $G_x$ by [8, p. 6-15, Cor. 2]. For the proof of (a) we may imbed $x$ in a
maximal torus $T$ and use the above notation. If $y$ in $G_x$ is written
$y=u\sigma_w b$ as in 2.4 then the uniqueness in 2.4 implies that $u,\sigma_w$ and $b$
are in $G_x$. By 2.1 and 2.2 we get:

2.8 \quad$G_x$ \uline{is generated by} $T$, \uline{those} $X_\alpha$ \uline{for which} $\alpha(x)=1$, \uline{and those}
$\sigma_w$ \uline{for which} $wx=x$. Then $G_{x0}$ is generated by $T$ and the $X_\alpha$ alone
because the group so generated is connected and of finite index in $G_x$ (see
[8, p. 3-01, Th. 1]). Let $G'$ be the group generated by the $X_\alpha$ alone,
and let $T'$ be the identity component of the intersection of the kernels of
the roots $\alpha$ such that $\alpha(x)=1$. Then $G'$ is semisimple by [8, p. 17-02,
Th. 1], and the other assertions of (a) are soon verified.

2.9 \quad Corollary. \uline{In} 2.7 \uline{every maximal torus containing} $\deleted{\unsure{x}}x$ \uline{also contains}
$T'$.

For in the above proof $T$ was chosen as an arbitrary torus containing $x$.

2.10 \quad\uline{Remark}. That $G_x$ in 2.7 need not be connected, even if $x$ is
regular, is shown by the example: $G=\mathrm{PSL}(2)$, $x=\operatorname{diag}(i,-i)$, $i^2=-1$.

% archive sheet number: 108

\scanpage{110}

\diagram{\begin{tikzcd}[ampersand replacement=\&,column sep=small,row sep=small]
\mathbb Q'(i)\arrow[r,dash]\arrow[d,dash] \& L\arrow[d, dash]\arrow[dr,dash] \& \deleted{\unsure{\mathbb Q'(i)}}\\
\mathbb Q'\arrow[r,dash] \& K\arrow[dr,dash]\arrow[dd,dash] \& L_0\arrow[d,dash] \&[-2ex] \bigg)\ \text{\unsure{galoisienne|galoisiens}}\\
\& \& K_0\\
\& \mathbb Q
\end{tikzcd}}

On veut construire le gpe de \unsure{t.m.} \inserted{$G(q)$} associé
à l'ens $M(q)$ des $\lambda\in L$ tels que $\lambda\bar\lambda=q^i$,
et $\lambda$ entiers exc. $p$. On a vu que

\[
M_0(q)/\mathrm{Torsion}\unsure{\subset|-|\simeq}(\mathbb Z^P)'\tag*{\unsure{(\ensuremath{*})|(\ensuremath{\pi})}}
\]

% slanted
\margin{N.B. $M_0(q)\inserted{=M_0(p)}$ \unsure{(ens des|= ens. des} $u\in L$ qui sont des unités exc. $p$ et satisfont $u\bar u=1$.}
$P$ l'ens des idéaux premiers de $L$ sur $p$.
On \unsure{peut supposer|doit supposer|l'a supposé} que les idéaux $\unsure{\mathfrak p|p|\text{pre}}$ de $K$ se \uline{décomposent} dans $L$.
Choisissons un des \deleted{\unsure{ces}} idéaux \unsure{$\mathfrak p$ au-dessus de $p$}, ce qui donne
une forme linéaire can. sur $M$

\[
\lambda\unsure{\longmapsto|\longrightarrow}v_{\mathfrak p}(\lambda),
\]

qui engendre d'ailleurs $\check M$ mod torsion sous \deleted{\unsure{$G$}} le
gpe de Galois $\Gamma$ de $L$, d'où

\[
{\mathbb G_m}_L\xrightarrow{v_L}G_L
\]

Soit $\Gamma_{\mathfrak p}$ le \unsure{grpe|gpe} de décomposition pour $\mathfrak p$, \deleted{qui est
\unsure{aussi}} et soit \deleted{$K_0$ le corps} $L_0$ le corps des
invariants de $\Gamma_{\mathfrak p}$, $K_0=L_0\cap K$. Alors
$L_0$ est le plus petit corps de rationalité de $v$,
qui s'écrit \unsure{donc|dans}

\[
v:{\mathbb G_m}_{L_0}\longrightarrow G_{L_0}
\]

D'autre part, si $\mathfrak p_0$ est induit par $\mathfrak p$ sur $L_0$,
alors \deleted{\unsure{tout} l'anneau de valuation complété} \inserted{l'anneau local complété}
de $\mathfrak p_0$ est \deleted{égal à celui de} isom : à celui de
$\mathbb Z_{\ell}$. Notons que le couple $(L_0,\mathfrak p_0)$ est défini
à isom unique près par la donnée de la \uline{classe}
de $L$, et peut se définir \deleted{\unsure{comme}} \inserted{via} \unsure{le|la} hensélis\unsure{é|e}
relatif de $\mathbb Z_{(p\mathbb Z)}$ dans $L$ \ldots. Noter que $L_0$ est
purement imaginaire, grâce à l'hypothèse
que $\tau\mathfrak p\ne\mathfrak p$.

% archive sheet number: 109

\scanpage{111}
\mspage{\handwritten{86}}

Utilizing (14) and assuming $n$ to be sufficiently
large, we find from (21) and (22):

\[
\handwritten{\dim H^1(V,O_V[nE'])=\dim H^1(V,O_V(-\mathscr E'_\nu)[nE']),}\tag*{\handwritten{(23)}}
\]

We can assume that $m'\geq n(\handwritten{\mathscr E'}_\nu)$ and so great that,
at $n=m'$, (23) is fulfilled. Then we have an exact
sequence

\diagram{\handwritten{\begin{tikzcd}[ampersand replacement=\&]
H^0(V,O_V[mE'])\arrow[r,"{r^*}"] \& H^0(\mathscr E'_\nu,O_{\mathscr E'_\nu}[m'E'])\arrow[r] \& 0
\end{tikzcd}}}

Utilizing the properties of the homomorphism
$r^*$ (See p. 208) we obtain hence that the mapping of the
variety $\handwritten{\mathscr E'}_\nu$ into a projective space, corresponding to
the group $H^{\unsure{0|O}}\handwritten{\big(}(\handwritten{\mathscr E'}_\nu),O_{\handwritten{\mathscr E'}_\nu}[m'E'])_{\handwritten{\mathscr E'}_\nu}\handwritten{\big)}$, coincides with the
restriction on $\handwritten{\mathscr E'}_\nu$ of the mapping $f_{m'}$. In particular,
we obtain that $f_{m'}$ maps $\handwritten{\mathscr E'}_\nu$ in a one-to-one manner
(that is, different points are transfered into different
points).

Let $q'$ be so great that $r=2^{q'}m-m'>0$. Let
us show that $f_{2^{q'}m}$ maps $\handwritten{\mathscr E'}_\nu$ in a one-to-one manner.
Let $P_1$ and $P_2$ be any two points on $\handwritten{\mathscr E'}_\nu$. It is
sufficient for us to show that there exists a cross-section of $H^0(V,O_V[2^{q'}mE'])$, which is equal to zero
in $P_1$, but not in $P_2$. It follows from above that there

% archive sheet number: 110

\scanpage{112}

\diagram{\begin{tikzcd}[ampersand replacement=\&,column sep=small,row sep=small]
\unsure{1|0}\arrow[r,dash] \& \mathbb Z/\unsure{n|2}\mathbb Z\arrow[r,dash]\arrow[d,dash] \& \Gamma\arrow[r]\arrow[d,phantom,"{\cup}" description] \& \Gamma'\arrow[r,dash] \& 0\\
\& 1\arrow[r,dash] \& \Gamma_{\mathfrak p}
\end{tikzcd}}

D'autre part, $(\mathbb Z^P)'$ correspond à
un tore sur $\mathbb Q$, qui peut se décrire ainsi, comme
on voit par \unsure{th.} de Galois : l'ext. quadratique
$L_0$ de $K_0$ définit une forme tordue
$U$ de $\unsure{\mathbb G_m}$ (savoir $\operatorname{Ker}\left(\prod_{L_0/K_0}{\mathbb G_m}_{L_0}\xrightarrow{\text{Norme}}{\mathbb G_m}_{K_0}\right)$\unsure{?})
et on considère \unsure{le}

\[
G'=\prod_{K_0/\mathbb Q}U=\operatorname{Ker}\left(\prod_{L_0/\mathbb Q}{\mathbb G_m}_{L_0}\unsure{-|\longrightarrow}\prod_{K_0/\mathbb Q}{\mathbb G_m}_{K_0}\right)
\]

(donc $G'(\mathbb Q)=\operatorname{Ker}\left(L_0^*\mathrel{\overset{N_{L_0/K_0}}{\rule{1.7em}{.4pt}}}K_0^*\right)$). On
a donc un hom

\[
\varphi:G'\longrightarrow G\tag*{(\ensuremath{**})}
\]

qui est une isogénie \uline{de $G'$ sur $G^{\unsure{ou|un|0u|\circ}}$}

D'ailleurs $v_{\unsure{L|L_0|\mathfrak p}}:{\mathbb G_m}_{L_0}\unsure{-|\longrightarrow}G_{L_0}$ se factorise
par $G'_{L_0}$

\diagram{\begin{tikzcd}[ampersand replacement=\&,row sep=small]
{\unsure{\text{soit}|\text{d'où}}\quad{\mathbb G_m}_{L_0}}\arrow[r,dash,"{v'}"]\arrow[d,phantom,"{\Vert}" description] \& G'_{L_0}\arrow[d]\\
{\mathbb G_m}_{L_0}\arrow[r,dash,"{v}"] \& G_{L_0}
\end{tikzcd}}

On fera attention que l'élémt $f_0$ \inserted{$=f\,j(p^{1/2})$} de $G^{ou}(\mathbb Q)\deleted{{}^{\unsure{*}}}$
ne se relève en un élémt \deleted{\unsure{de}} $f'$ de $G'(\mathbb Q)$ que
si \deleted{\unsure{$\mathbb Q$|$\varphi$}}\inserted{$\varphi$} est un isom i.e. $M(q)/\mathrm{torsion}\xrightarrow{\sim}(\mathbb Z^P)'$.

% archive sheet number: 111

\scanpage{113}
\mspage{\handwritten{84}}

Since $Q\in(x-C_\beta)_{\unsure{0|O}V'}\cap  U'$, then $Q\in h((x-C_\beta)_{0V^*}\cap  U^*$,
for all $C_\beta$. But the sets $(x-C_{\beta'})_{0V^*}\cap  U^*$ and
$(x-C_{\beta''})_{0V^*}\cap  U^*$ do not intersect (the (point) $x$ is regular
in $U^*$ and can not take two different values $C_{\beta'}$ and $C_{\beta''}$
in a point of $U^*$. Hence it follows that $Q$ has infinitely
many preimages on $U^*$, that is, on $V^{*(2^q m)}$.

We arrived at a contradiction. Consequently, $O_{Q'}=O_Q$.
If $Q''$ is another point on $V^{*(2^q m)}$ such that $hQ''=Q$,
then, from $O_{Q''}=O_Q$, we obtain $O_{Q'}=O_{Q''}$ and $Q'=Q''$.
Namely, $h$ is a one\deleted{\unsure{=}}-to-one mapping. Hence it follows that
$h_{q^q m,2^{q_{\unsure{o|0}}}m}$ is a one-to-one mapping on $V^{'(2^q m)}$.

In a projective space, connected with the mapping $f_{2^{q_o}m}$,
let us examine a hypersurface, containing $W$ and not
containing $V^{(2^{q_o}m)}$.

Evidently, the cross-section of the variety $V^{(2^{q_o}m)}$
with this hypersurface is the image by $f_{2^{q_o}m}$ of the
subvariety $\handwritten{\mathscr E'}_\nu$ of the variety $V$, which is the variety of
\margin{p. 213}zeros of some cross-section of the group $H^{\unsure{0|O}}(V,O_V[\nu  E'])$,
where $[\nu  E']$ is some fibre space of straight lines, a
multiple by the fibre $[2^{q_o}mE']$. According to the induction
assumption, there exists such $n(\handwritten{\mathscr E'}_\nu)$, that, at $n\geq n(\handwritten{\mathscr E'}_\nu)$

\[
\handwritten{H^1\left(\mathscr E'_\nu,O_{\mathscr E'_\nu}[nE']_{\mathscr E'_\nu}\right)=0}
\]

and such $m'$, that a mapping of the variety $\handwritten{\mathscr E'}$ into a
projective space, corresponding to the group $H^0(\handwritten{\mathscr E'}_\nu,O_{\handwritten{\mathscr E'}_\nu}[m'E']_{\handwritten{\mathscr E'}_\nu})$
becomes a biregular imbedding. Let us examine the exact

% archive sheet number: 112

\scanpage{114}

% slanted
\margin{demander à Serre ?}
Ce n'est en général pas le cas, \unsure{$=$|même} si
$K_{\unsure{8|0}}=\mathbb Q$ : il suffit de donner un exemple où
$\mathbb Z/2\mathbb Z$ opère non \unsure{trivialmt|trivialement} sur $\mathrm{Pic}(A)$, $A$ étant l'\unsure{anneau}
\deleted{\unsure{groupe}} des entiers de $L$. On constate d'ailleurs que
$\mathbb Z/2\mathbb Z$ opère sur $\mathrm{Pic}(A)$ par symétrie ; il suffit donc
de trouver $\mathrm{Pic}(A)$ non un $2$-gpe !.

% archive sheet number: 113

\scanpage{115}
\mspage{\handwritten{85}}

sequence

\diagram{\handwritten{\begin{tikzcd}[ampersand replacement=\&,column sep=small]
0\arrow[r] \& O_V(-\nu  E')\arrow[r] \& O_V(-\mathscr E'_\nu)\arrow[r] \& G\arrow[r] \& 0\unsure{,|.}
\end{tikzcd}}}

Here $G$ is some coherent sheaf, concentrated on $\handwritten{\mathscr E'}_\nu$.
It follows from the observation made in the beginning of
the proof that, beginning with some sufficiently large $n$,

\[
\handwritten{H^1(V,G[nE'])=0.}
\]

Hence it follows that, beginning with some sufficiently
large $n$,

\[
\handwritten{\dim H^1(V,O_V(-\mathscr E'_\nu)[nE'])\leq\dim H^1(V,O_V[(n-\nu)E']).}\tag*{\handwritten{(21)}}
\]

or

\[
\handwritten{\dim H^1(V,O_V(-\mathscr E'_\nu)[nE'])\leq\dim H^1(V,O_V[(n-\nu)E']).}
\]

From the exact sequence

\diagram{\handwritten{\begin{tikzcd}[ampersand replacement=\&,column sep=small]
0\arrow[r] \& O_V(-\mathscr E'_\nu)\arrow[r] \& O_V\arrow[r] \& O_{\mathscr E'}\arrow[r] \& 0
\end{tikzcd}}}

we obtain, at $n\geq n(\handwritten{\mathscr E'}_\nu)$:

\[
\handwritten{\dim H^1(V,O_V[nE'])\unsure{\leq|\leqq}\dim H^1(V,O_V(-\mathscr E'_\nu)[nE']),}\tag*{\handwritten{(22)}}
\]

% archive sheet number: 114

\scanpage{117}

Soit $G$ un grpe \deleted{\unsure{pro}} alg comm sur $\mathbb Q$ provenant
d'une situation motivique \inserted{$\unsure{M|\mathcal M}$}, soit $U=\operatorname{Ker}(G^0\mathrel{\mathop{\unsure{-|\longrightarrow}}\limits^{\epsilon}}\mathbb G_m)$.
On a \unsure{donc|encore} un élément canonique dans $\mathbb G_m$.
\deleted{On a \unsure{$U$ tore|$U=$} \unsure{ssi|$G$}.}
% slanted
\margin{N.B. \unsure{Je n'utilise|La condition} que les \unsure{élém.|éléments|idées} de structure $1,2,3$}

\uline{Proposition} 1. Conditions équivalentes

(i) $U$ est un tore

(ii) $\epsilon$ est primitif dans l'ens des caractères $G^0\mathrel{-}\mathbb G_m$

(iii) $M_{\mathbb R}$ admet un foncteur fibre \inserted{``normalisé''}, i.e. $\unsure{\mathscr G|\mathcal G}_{\mathbb R}\ne\varnothing$.
% slanted
\margin{i.e. compatible \unsure{sur les|avec les|à un} motifs eff. de poids $0$ avec le foncteur fibre standard \unsure{corresp.} \unsure{canonique}.}

C'est tjrs vrai si $M$ a degrés pairs, i.e. $j(\mu_2)=0$.

\uline{Proposition} 2. Conditions équivalentes

(i) $U$ est de la forme $U^0\times\mu_2$, $U^0$ un tore

(ii) $\epsilon$ est \deleted{\unsure{un prim.}} un carré \inserted{dans $\operatorname{Hom}(G^0,\mathbb G_m)$}

(iii) $M_{\mathbb R}$ n'a pas de foncteur fibre \inserted{``normalisé''}, i.e. $\mathscr G_{\mathbb R}=\varnothing$.

Le fait que (i) équivaut à (ii) \inserted{dans prop 1 resp. 2}
résulte trivialmt de \deleted{l'\unsure{équivalence} de \unsure{ce}} l'existence de $j$ avec $\epsilon j(\lambda)=\lambda^2$,
ainsi que le fait que $(1\,\mathrm{ii})\Longleftrightarrow(\text{non }\deleted{\unsure{\mathrm{(iii)}}}\,2\,\mathrm{(ii)})$.

On a $(\mathrm{iii})\Rightarrow(\mathrm{ii})$, car plus généralement pour tt
\deleted{\unsure{fonct.}} \inserted{objet} de \deleted{\unsure{degré} \unsure{objet|degré}} \inserted{degré} impair de $M_{\mathbb R}$, le \deleted{degré}\inserted{rang} doit
être pair, comme il résulte de l'existence d'une forme alternée fondamentale. D'autre
part $(\mathrm i)\Rightarrow(\mathrm{iii})$, car si $T$ est un tore
compact sur $\mathbb R$, on a $H^{\unsure{2|1|n}}(\mathbb R,T)=0$ ! Enfin, si
$M$ est à degrés pairs i.e. $j(\mu_2)=0$, alors $j=2j'$
\deleted{\uline{Cor}} et $\epsilon j'=\mathrm{id}$, donc on a (ii).

\uline{Cor} Si $M$ est quelconque, la sous-catégorie \inserted{motivique} $M^{\mathrm{pair}}$ \deleted{\unsure{admet}}
(à qui correspond $G/\deleted{\unsure{T}}j(\mu_2)$) \deleted{\unsure{est}} admet un
$\mathbb R$-foncteur fibre.

% archive sheet number: 116

\scanpage{118}

% Two long oblique pen strokes traverse the typed draft.
\deleted{Page 33. Ligne $-9$ à $-1$, remplacer le texte par le suivant:}

\deleted{\handwritten{«} \uline{Remarque}s7.10. a) Lorsque $G$ est à fibres connexes, alors on peut
dans c) abandonner l'hypothèse que les $\deleted{\unsure{p}}\,{}_{p^{\handwritten{\unsure{\nu|v}}}}H'$ soient séparés. En
effet, avec les réductions faites dans la démonstration, on peut
supposer $u$ bijectif, donc $H$ à fibres connexes, donc $H$ séparé en vertu du théorème de séparation de Raynaud (Exp $\mathrm{VI}_{\mathrm B}$).}

\deleted{b) D'autre part, des exemples de Raynaud montrent qu'on
ne peut en général dans c) abandonner ni l'hypothèse que les $ {}_{p^{\handwritten{\nu}}}G'$
soient finis sur $S'$, ni que les $ {}_{p^{\handwritten{\nu}}}H'$ soient séparés sur $S'$.}

\deleted{c) \inserted{On peut montrer que} \deleted{L}la conclusion de a) est encore valable si, au lieu de supposer $G$ à fibres connexes, on le suppose de type multiplicatif; \deleted{\unsure{i.e.}}
tout monomorphisme $u:G\handwritten{\longrightarrow}H$ de $S$-préschémas en groupes\deleted{\unsure{,}} de présentation finie, avec $G$ de type multiplicatif, \deleted{\unsure{et H|et N}} est une immersion. \handwritten{»}}

\deleted{Ligne 14, lire $S'$ au lieu de $S$.}

\handwritten{Page 35, ligne 9,10, supprimer ``lisse sur $k$''.}

% archive sheet number: 117

\scanpage{119}

\uline{Cor} Dans le cas de la prop 2, \deleted{la catégorie} on a
$G^0\simeq U^0\times\mathbb G_m$, $\epsilon=\mathrm{pr}_2$, $j(\lambda)=\deleted{\unsure{\mathrm{pr}(\lambda)^2|\mathrm{inj}(\lambda)^2}}\,\mathrm{inj}_{1}(\lambda)^2$. La structure
\deleted{motivique $1)\,2)\,3)$} de catégorie tensorielle de $\deleted{\unsure{G}}M_{\mathbb R}$ est
connue totalement par la connaissance de
$G$ et l'élémt de $H^{\unsure{2|1}}(\mathbb R,\deleted{\unsure{G}}U)$ caractéristique, qui \deleted{\unsure{puis.}}
\deleted{est} l'unique élémt non nul de $H^2(\mathbb R,U)$,
correspondant à l'unique élément non nul
de $H^2(\mathbb R,\mu_2)$. Il existe un \deleted{\unsure{élém}} objet $E$ de
$M_{\mathbb R}$ \deleted{\unsure{en}} de degré 1, rang 2, \deleted{sur lequel} tel que
la sous-catégorie tensorielle $\deleted{\unsure{M'}}M'$ qu'il engendre
soit déjà $\mathbb R$-tordue, et \deleted{\unsure{correspond} \unsure{au} \unsure{fait que}}
\deleted{\unsure{sur}} que $U^0$ ``y \unsure{opère|agisse} trivialmt''.
\par\hspace*{.55\linewidth}\rule{2em}{.3pt}

Dans le cas favorable $\unsure{\mathscr G|\mathcal G}_{\mathbb R}\ne\varnothing$, \deleted{\unsure{on a en gros}} \inserted{i.e. $U$ un tore}
l'ens $\deleted{\operatorname{Ob}\mathscr G_{\mathbb R}}$ \inserted{$H^0(\mathscr G_{\mathbb R})$ des objets à iso. près de $\mathscr G_{\mathbb R}$} est un
\deleted{\unsure{ens.} principal homogène} \inserted{torseur}
sous le grpe $H^{1}(\mathbb R,U)$, qui lui-$=$ est
isom. can à $ {}_2U(\mathbb R)\underset{\substack{\mathrm{isom}\\\mathrm{can.}}}{\simeq}(\mathbb Z/2\mathbb Z)^{\dim U}$. D'autre
part, à chaque élémt de $\operatorname{Ob}\mathscr G_{\mathbb R}$ correspond
une \deleted{structure de \unsure{positivité}} \inserted{structure polarisante} sur $G^0$,
définie par un hom $\mu_{2_{\mathbb C}}\mathrel{\mathop{\longrightarrow}\limits^{i_1}}G^0_{\mathbb C}$ tel que
$\epsilon_{\mathbb R}i_1=\mathrm{id}$, $i_{\unsure{1|2}}\bar i_1=j_{\mathbb R}$, et l'ens de ces structures est
égalmt un torseur sous $ {}_2U(\mathbb R)$. L'application

\[
H^0(\mathscr G_{\mathbb R})\longrightarrow\operatorname{Pol}(\mathscr G_{\mathbb R})
\]

est compatible avec \deleted{\unsure{le}} l'opération de $ {}_2U(\mathbb R)$, \uline{donc}
\uline{bijective}. Donc les $\mathbb R$-foncteurs fibres à \unsure{isom|iso|isom.}

% archive sheet number: 118
